% Weekly exercise sheet for Lecture 14; also included in the synthesis sheet.
\ifdefined\ExercisesSynthesis
\else
\documentclass[11pt]{article}
\usepackage[utf8]{inputenc}
\usepackage[T1]{fontenc}
\usepackage[english]{babel}
\usepackage{lmodern,microtype}
\usepackage[margin=0.9in]{geometry}
\usepackage{amsmath,amssymb,amsthm,mathtools}
\usepackage{enumitem,hyperref}
\newcommand{\E}{\mathbb E}
\newcommand{\Qbb}{\mathbb Q}
\newcommand{\dd}{\,\mathrm d}
\newcommand{\F}{\mathcal F}
\theoremstyle{definition}
\newtheorem{exercise}{Exercise}[section]
\newlist{exparts}{enumerate}{1}
\setlist[exparts,1]{label=\textbf{(\alph*)},itemsep=0.3em,topsep=0.4em,leftmargin=2em}
\newcommand{\lecturesummary}[1]{\begin{quote}\small\textbf{Focus.} #1\end{quote}}
\title{Stochastic Calculus and Option Pricing\\Lecture 14 Exercises}
\author{}
\date{}
\begin{document}
\maketitle
\vspace{-1.2em}
\setcounter{section}{13}
\fi

\section{Week 14 --- Pricing Toolbox: A Guided Refresher}
\lecturesummary{A final theoretical variation on the product rule for
stochastic processes.}

\begin{exercise}[Quadratic covariation controls terminal covariance]
Let \(X\) and \(Y\) be bounded continuous martingales on \([0,T]\), with
deterministic initial values.
\begin{exparts}
  \item Apply the product rule to \(X_tY_t\), keeping the quadratic
  covariation term explicit.
  \item If \(\E[\langle X,Y\rangle_T]=0\), take expectations and show that
  \(\operatorname{Cov}(X_T,Y_T)=0\).
  \item Give a simple example of dependent, uncorrelated random variables
  to explain why zero terminal covariance does not imply independence.
\end{exparts}
\end{exercise}

\ifdefined\ExercisesSynthesis\else
\end{document}
\fi
